% concatenated from main__3_.tex

%% Elsevier CAS single-column manuscript generated from the project source.
\documentclass[a4paper,fleqn]{cas-sc}

\usepackage[authoryear,longnamesfirst]{natbib}
\usepackage{amsmath,amssymb,amsfonts,amsthm,mathtools}
\usepackage{enumitem}
\usepackage{booktabs}

\newtheorem{theorem}{Theorem}[section]
\newtheorem{corollary}[theorem]{Corollary}
\newtheorem{proposition}[theorem]{Proposition}
\newtheorem{lemma}[theorem]{Lemma}
\newtheorem{conjecture}[theorem]{Conjecture}
\newtheorem{problem}[theorem]{Problem}
\theoremstyle{definition}
\newtheorem{definition}[theorem]{Definition}
\newtheorem{example}[theorem]{Example}
\theoremstyle{remark}
\newtheorem{remark}[theorem]{Remark}

\newcommand{\R}{\mathbb R}
\newcommand{\Q}{\mathbb Q}
\newcommand{\Z}{\mathbb Z}
\newcommand{\F}{\mathbb F}
\newcommand{\N}{\mathbb N}
\newcommand{\one}{\mathbf 1}
\newcommand{\Aalg}{\mathcal A}
\newcommand{\Walg}{\mathcal W}
\newcommand{\Aut}{\operatorname{Aut}}
\newcommand{\Stab}{\operatorname{Stab}}
\newcommand{\Orb}{\operatorname{Orb}}
\newcommand{\End}{\operatorname{End}}
\newcommand{\diag}{\operatorname{diag}}
\newcommand{\tr}{\operatorname{tr}}
\newcommand{\rank}{\operatorname{rank}}
\newcommand{\dist}{\operatorname{dist}}
\newcommand{\diam}{\operatorname{diam}}
\newcommand{\Span}{\operatorname{span}}
\newcommand{\Sel}{\operatorname{Sel}}
\newcommand{\M}{M}
\newcommand{\U}{U}
\newcommand{\PiAD}{\Pi_{\!AD}}
\newcommand{\Had}{\circ}
\newcommand{\homG}{\operatorname{hom}}

% The CAS class prints an empty ORCID footnote when no ORCID is supplied.
% Keep the footnote only when an ORCID has actually been declared.
\ExplSyntaxOn
\RenewDocumentCommand \printorcid { }
 {
  \seq_if_empty:NF \g_stm_orcid_seq
   {
    \group_begin:
     \tex_let:D \thefootnote \relax \footnotetext
      {
       \raggedright
       \textsc{orcid}(s):\c_space_token
       \seq_use:Nn \g_stm_orcid_seq { ;~ }
      }
    \group_end:
   }
 }
\ExplSyntaxOff

\begin{document}
\let\WriteBookmarks\relax
\def\floatpagepagefraction{1}
\def\textpagefraction{.001}

\shorttitle{Principal adjacency-degree moments}
\shortauthors{Z. Lu}

\title [mode = title]{Adjacency-degree algebras and spectral determination of graphs}

\author[1]{Zhipeng Lu}
\cormark[1]
\ead{zhipeng.lu@hotmail.com}

\affiliation[1]{organization={Shenzhen MSU--BIT University and Guangdong Laboratory of Machine Perception and Intelligent Computing},
            city={Shenzhen},
            postcode={518172},
            state={Guangdong},
            country={China}}

\cortext[1]{Corresponding author}

\author[2]{Pengxiang Li}
\ead{pengxiangli@bit.edu.cn}

\affiliation[2]{organization={Beijing Institute of Technology},
            city={Beijing},
            postcode={100081},
            state={Beijing},
            country={China}}


\begin{abstract}
McKay proved that the spectra of all polynomial functions of the adjacency
matrix $A$ and the diagonal degree matrix $D$ determine a tree.  We prove a
principal version of this theorem.  Let $\mathcal A(G)=\langle I,A_G,D_G\rangle$
and let $M_G=\mathcal A(G)\mathbf1$ be the cyclic module generated by the
all-ones vector.  For connected graphs the ideal $\mathcal A(G)J\mathcal A(G)$,
where $J=\mathbf1\mathbf1^T$, acts on $M_G$ as the full endomorphism algebra.  We
show that every forest satisfies $M_G=U_G$, the automorphism-orbit module, and
that the induced algebra on the orbit quotient of a tree is a full matrix algebra.
It follows that the scalar moments $\mathbf1^Tw(A_T,D_T)\mathbf1$ determine every
tree.  For general graphs these moments are degree-decorated caterpillar
homomorphism counts.  The resulting moment-rigidity class lies inside the
amenable, compact, refinable hierarchy of color refinement, and its first
small-order failures are ten-vertex integral switchings invisible to $M_G$.
\end{abstract}

\begin{highlights}
\item The principal ideal acts as the full endomorphism algebra of $M_G$.
\item Forests are AD-observable, and trees are determined by principal AD moments.
\item Principal AD moments are degree-decorated caterpillar counts.
\end{highlights}

\begin{keywords}
Spectral graph theory \sep adjacency-degree algebra \sep tree reconstruction \sep color refinement \sep caterpillar homomorphisms
\end{keywords}

\maketitle

\section{Introduction}
\label{sec:introduction}

Spectral graph theory asks how much of a graph is encoded by linear algebraic
invariants.  Since the early examples of \citet{CollatzSinogowitz1957}, it has been clear
that the adjacency spectrum alone is usually not a complete invariant.
\citet{Schwenk1973} proved that almost all trees have a cospectral mate, and
Godsil--McKay switching gives a general construction of cospectral
graphs~\citep{GodsilMcKay1982}.  Surveys of graphs determined by their spectra
appear in~\citep{vanDamHaemers2003,vanDamHaemers2009}, and the standard
algebraic background may be found in~\citep{CvetkovicDoobSachs1980,GodsilRoyle2001}.

McKay's theorem is a striking positive result in this landscape: every tree is
determined by the spectra of all polynomial functions of its adjacency matrix
$A$ and diagonal degree matrix $D$~\citep{McKay1977}.  Equivalently, the pair
$(A,D)$ contains enough spectral information to reconstruct a tree. In this work, we consider a sharper form of this statement. Let
\[
        \mathcal A(G)=\langle I,A_G,D_G\rangle
\]
be the adjacency-degree algebra, and let
\[
        M_G=\mathcal A(G)\mathbf1
\]
be the cyclic module generated by the all-ones vector.  We study the scalar
principal moments
\[
        \mu_G(w)=\mathbf1^T w(A_G,D_G)\mathbf1,
        \qquad w\in\mathbb R\langle x,y\rangle .
\]
The main theorem says that these moments determine every tree.

The proof is deliberately close to McKay's leaf-stripping method, and we view it
as a linear-algebraic implementation of that method inside the principal cyclic
module.  The single matrix identity driving the argument is the following: if
$W$ is a set of leaves of a forest and $X_W$ is the diagonal indicator of $W$,
then
\[
        A X_W A
\]
is diagonal and its $v$-th diagonal entry is the number of selected leaves
adjacent to $v$.  This identity is the matrix form of the leaf-neighbor counts
used in McKay's pruning.  Its significance here is that the required selectors
are produced by ordinary multiplication in the algebra generated by $A$ and $D$.
Thus the pruning process yields orbit indicators from the principal module
rather than from the full adjacency-degree spectral calculus.

The algebraic framework is as follows.  Because $A$ and $D$ are symmetric,
$\mathcal A(G)$ is a transpose-closed finite-dimensional real matrix algebra,
and hence semisimple.  We do not classify all of its simple summands.  Instead,
we isolate the canonical summand selected by $\mathbf1$.  For connected graphs,
$J/n$, where $J=\mathbf1\mathbf1^T$, is the Laplacian zero-eigenspace projection,
so $J\in\mathcal A(G)$.  We prove
\[
        \mathcal A(G)J\mathcal A(G)
        =\operatorname{span}\{uv^T:u,v\in M_G\},
\]
and this ideal acts on $M_G$ as $\operatorname{End}(M_G)$.

The graph-theoretic content is measured by the modules
\[
        M_G\subseteq W_G\subseteq U_G,
\]
where $W_G$ is the stable color-refinement module and $U_G$ is the automorphism-orbit
module.  The equality $W_G=U_G$ is the standard refinability condition of
\citet{ArvindKoblerRattanVerbitsky2017}; in their hierarchy compact graphs are
refinable, and amenable graphs are compact.  The graphs determined by the
principal moments therefore sit in the established chain
\[
        \{\text{graphs determined by principal moments}\}
        \subseteq \{\text{amenable graphs}\}
        \subseteq \{\text{compact graphs}\}
        \subseteq \{\text{refinable graphs}\}.
\]
Note that compactness in Tinhofer's sense is integrality of the fractional automorphism polytope;
compact graphs are refinable and amenable graphs are compact
\citep{Tinhofer1991,ArvindKoblerRattanVerbitsky2017}. Thus compactness explains
one standard route to the right-hand equality $W_G=U_G$, while our work studies
the left-hand principal inclusion $M_G\subseteq W_G$ and the scalar moments of
$M_G$.  This inclusion is generically an equality for the elementary reason that
if the ordinary walk matrix
\[
        [\one,A\one,A^2\one,\ldots,A^{n-1}\one]
\]
has full rank, then its columns already span $M_G=\mathbb R^{V(G)}$, while almost all graphs are controllable in this sense by the theorem of \citet{ORourkeTouri2016}.
The right-hand equality alone, however, does not force the left-hand equality.
For example, the 8-vertex graph \verb|G?qa`g| is refinable: its stable
color-refinement partition agrees with its automorphism-orbit partition and has
five cells.  Nevertheless its principal cyclic module has dimension four, so
$M_G\subsetneq W_G=U_G$.  This density-one observation is only module-theoretic.
The forest theorem proved below gives a structured, non-generic family with
$M_F=U_F$ for every forest $F$.  For trees, the orbit quotient carries more structure:
the operators induced by $A$ and $D$ on the quotient generate a full matrix algebra.
Consequently the principal moments recover the weighted orbit quotient, and the
weighted orbit quotient recovers the tree.

For comparison, one-parameter spectral refinements are not enough for trees. The
generalized characteristic polynomial of Wang--Li--Lu--Xu,
\[\det(\lambda I-(A-tD)),\]
is a natural slice of the adjacency-degree calculus~\citep{WangLiLuXu2011}.  It is
well known from McKay's construction that such one-parameter data can coincide
for non-isomorphic trees.  We include a small explicit pair in
Section~\ref{sec:tree-reconstruction}.  Related recent work on degree-similar
graphs studies the unpointed simultaneous-similarity problem for the pair
$(A,D)$. This is distinct from the pointed cyclic moment problem considered
here, since simultaneous similarity preserves trace words whereas principal
moments depend on the distinguished vector $\one$~\citep{GodsilSun2025DegreeSimilar,FanXingZhangWang2025DegreeSimilar}.

For general graphs, the same principal moments have a homomorphism-count
interpretation.  The monomial moment
\[
        \mathbf1^TD^{a_0}AD^{a_1}A\cdots AD^{a_\ell}\mathbf1
\]
is the homomorphism count from the \emph{caterpillar} obtained from a path by attaching
$a_i$ leaves to the $i$-th spine vertex.  Hence principal AD-rigidity (section~\ref{sec:caterpillar-rigidity}) is a rigidity
under degree-decorated caterpillar counts.  By the theorem of Dell--Grohe--Rattan,
color refinement is equivalence under all tree homomorphism counts~\citep{DellGroheRattan2018}.
Thus principal AD-rigid graphs are amenable to color refinement. Complete
enumeration shows that the two equivalence relations coincide on graphs with at
most nine vertices, but ten-vertex integral switchings invisible to the
principal module give the first strict separation. These switchings explain the
gap as a structural one: principal moments record one-spine degree-decorated tree
data, while color refinement records full rooted branching data.

We organize the paper as follows.  Section~\ref{sec:algebra} records the
adjacency-degree algebra, the principal ideal, and the modules $M_G$, $W_G$, and
$U_G$.  Section~\ref{sec:trees} proves the forest orbit-module theorem.
Section~\ref{sec:tree-reconstruction} proves tree reconstruction from principal
moments and places it next to the Wang pencil.  Section~\ref{sec:algorithmic-extraction}
extracts finite certificates and a linear-time tree canonization from the
constructive proof.  Section~\ref{sec:caterpillar-rigidity} discusses principal
AD-rigidity, the comparison with color refinement, and the first integral
switching obstructions.

\section{The adjacency-degree algebra and its principal block}
\label{sec:algebra}

Let $G$ be a finite simple graph with vertex set $V$, adjacency matrix $A$, degree
matrix $D$, and all-ones vector $\one$.  We write
\[
        \Aalg(G)=\langle I,A,D\rangle\subseteq M_V(\R)
\]
for the real algebra generated by $I,A,D$.  It is closed under transpose (hence a *-algebra), which has the following consequence: $\Aalg(G)$ is semisimple as a real matrix
algebra.  Indeed, if $R$ is its Jacobson radical, then transpose maps $R$ to
itself, because transpose is an anti-automorphism and the radical is invariant
under anti-isomorphism.  Thus, for $X\in R$, also $X^T\in R$ and hence
$XX^T\in R$.  Since the radical is nilpotent, $XX^T$ is nilpotent; being symmetric,
it has only the eigenvalue $0$, so $0=\tr(XX^T)=\|X\|_F^2$ and $X=0$.  Hence the
radical is zero. While classifying simple
summands of $\Aalg(G)$ might be of independent interest, we focus on understanding the cyclic summand generated by $\one$:
\[
        M_G=\Aalg(G)\one.
\]

Let $\mathcal P_{\rm WL}(G)$ be the stable color-refinement partition, and put
\[
        W_G=\Span\{\one_C:C\in\mathcal P_{\rm WL}(G)\}.
\]
The stable partition is equitable and degree-homogeneous, so $W_G$ is invariant
under $A$ and $D$ and contains $\one$.  Since color classes are unions of
automorphism orbits,
\[
        M_G\subseteq W_G\subseteq U_G,
        \qquad
        U_G=\Span\{\one_O:O\in\Orb(\Aut(G))\}.
\]
We say that $G$ is \emph{principal color-observable} if $M_G=W_G$.  The
equality $W_G=U_G$ is the standard \emph{refinable} property: color refinement
produces exactly the automorphism orbit partition~\cite{ArvindKoblerRattanVerbitsky2017}.
Finally, $G$ is \emph{AD-observable} if $M_G=U_G$. Thus, AD-observability is the
conjunction of principal color-observability and refinability.  %These are module properties inside a fixed graph; they should not be confused with graph determination from the scalar moments $\mu_G$.

The inclusion $M_G\subseteq U_G$ is immediate: every automorphism of $G$ commutes
with $A$ and $D$ and fixes $\one$, so every vector in $M_G$ is constant on
automorphism orbits.  Equivalently, if $S$ is the orbit-indicator matrix and
$B,\Delta$ are the orbit quotient operators, then
\[
        M_G=S(\langle I,B,\Delta\rangle\one).
\]
Hence $G$ is AD-observable precisely when $\langle I,B,\Delta\rangle\one$ is the
full quotient space.

For connected graphs the vector $\one$ also selects a canonical ideal in the full
algebra.  Since the Laplacian $L=D-A$ has $\ker L=\R\one$, the orthogonal
projection $J/n$ onto this kernel is a polynomial in $L$ and belongs to
$\Aalg(G)$.

\begin{theorem}[The principal rank-one ideal]
\label{thm:principal-ideal}
Let $G$ be connected and put $J=\one\one^T$.  Then
\[
        \Aalg(G)J\Aalg(G)=\Span\{uv^T:u,v\in M_G\}.
\]
Moreover, after restriction to $M_G$, this ideal is the full endomorphism algebra
$\End(M_G)$.
\end{theorem}

\begin{proof}
Every element of $\Aalg(G)J\Aalg(G)$ is a linear combination of matrices
$p(A,D)Jq(A,D)$.  Since $J=\one\one^T$,
\[
        p(A,D)Jq(A,D)=(p(A,D)\one)(q(A,D)^T\one)^T .
\]
The algebra is transpose-closed, so both $p(A,D)\one$ and $q(A,D)^T\one$ belong
to $M_G$.  Hence the ideal is contained in the displayed rank-one span.
Conversely, every $u,v\in M_G$ has the form $u=p(A,D)\one$, $v=q(A,D)\one$, and
\[
        uv^T=p(A,D)Jq(A,D)^T
\]
lies in the ideal.  The rank-one operators $uv^T$ span all endomorphisms of
$M_G$ and vanish on $M_G^\perp$.
\end{proof}


\iffalse
Let $P_M$ be the orthogonal projection onto $M_G$.  Since $M_G$ is invariant
under the symmetric generators $A,D$, the orthogonal complement $M_G^\perp$ is
also invariant.  Thus $P_M$ commutes with $\Aalg(G)$, and Theorem~\ref{thm:principal-ideal}
identifies the principal central summand:
\[
        \Aalg(G)=P_M\Aalg(G)P_M\oplus (I-P_M)\Aalg(G)(I-P_M),
        \qquad
        P_M\Aalg(G)P_M=\End(M_G).
\]

\begin{remark}[The full trace character]
The full AD trace character
\[
        \tau_G(w)=\tr w(A_G,D_G)
\]
is stronger than the principal character on connected graphs.  Indeed, if $G$ is
connected, then $J/n=q(D-A)$ for a polynomial $q$ which is $1$ at the Laplacian
zero eigenvalue and $0$ at the other Laplacian eigenvalues. Therefore,
\[
        \mu_G(w)=\tr(Jw(A,D))=n\,\tr(q(D-A)w(A,D))
\]
is a linear combination of full trace moments. Thus, equality of all full AD
trace moments for two connected graphs implies equality of their principal AD
moments. We focus on the principal character because it is the
canonical summand selected by $\one$, because it already determines all trees,
and because it has the decorated-caterpillar interpretation of
Section~\ref{sec:caterpillar-rigidity}.
\end{remark}
\fi 



\paragraph{The full trace character vs principal character.} The full AD trace character$\tau_G(w)=\tr w(A_G,D_G)$ is strictly stronger than the principal character in general. The simplest examples are regular graphs whose principal character degenerates trivially. A more sophisticate non-regular example is as follows. 

\begin{example}[A 10-vertex principal-invisible switch]
\label{ex:n10-counterexample}

Let \(G_{10}\) be
the graph on \(\{0,\ldots,9\}\) with
\[
\begin{aligned}
E(G_{10})=\{&04,06,08,15,17,18,24,25,29,36,37,39,\\
            &47,48,56,59,67,68,79,89\},
\end{aligned}
\]
and \(H_{10}\) be obtained from \(G_{10}\) by the switch
\[
        04,15\quad\longleftrightarrow\quad 05,14 .
\]
The two graphs have the same degree sequence
\[
        (3,3,3,3,4,4,5,5,5,5).
\]
For \(G_{10}\), the principal module has dimension \(5\).  A certificate,
recorded in Appendix~\ref{app:n10-certificate}, gives a basis matrix \(K\) for
\(M_{G_{10}}\) with
\[
        A_{G_{10}}K=KC_A,\qquad D_{G_{10}}K=KC_D,
\]
and the switch perturbation \(E=A_{H_{10}}-A_{G_{10}}\) satisfies
\[
        EK=0.
\]
Thus \(D_{H_{10}}=D_{G_{10}}\) and \(A_{H_{10}}\) agrees with \(A_{G_{10}}\) on
the whole principal module.  Hence
\[
        \mathbf1^T w(A_{G_{10}},D_{G_{10}})\mathbf1
        =
        \mathbf1^T w(A_{H_{10}},D_{H_{10}})\mathbf1
        \qquad\text{for all words }w .
\]
The full trace character sees the switch immediately:
\[
        \operatorname{tr}(A_{G_{10}}^3)=30,
        \qquad
        \operatorname{tr}(A_{H_{10}}^3)=42 .
\]
So the full character distinguishes this principal mate, while the principal
character does not.  
\end{example}

It turns out that, for trees, this loss
of the non-principal summands does not occur: the principal character already
contains enough information to reconstruct the tree. This will be the focus of the subsequent sections.




\section{A matrix form of McKay leaf stripping}
\label{sec:trees}

This section proves the forest orbit-module theorem.  The construction follows
the leaf-stripping scheme used by McKay in the spectral characterization of
trees~\cite{McKay1977}.  The new point is that the peeling data are produced
inside the principal adjacency-degree module.  In particular, the leaf-neighbor
counts used by the pruning are realized by ordinary matrix products rather than
by adding external color selectors.

The proof has four steps.  First, the leaf diagonal identity converts selected
leaf sets into diagonal count operators.  Second, guarded pruning uses these
diagonal count operators to construct idempotents for lower rooted types.
Third, center selectors produce the seed orbit indicators at the terminal core.
Finally, reverse propagation moves these indicators back through the pruning
process.  The rooted orbit classification at the end verifies that the
constructed vectors are exactly the automorphism-orbit indicators.

\subsection{The leaf diagonal identity}

Let \(F\) be a forest on vertex set \(V\), with adjacency matrix \(A_F\).  For
\(W\subseteq V\), write \(X_W=\diag(\one_W)\).

\begin{lemma}[Leaf diagonal lemma]
\label{lem:leaf-diagonal}
If \(W\) is a set of leaves of \(F\), then \(A_FX_WA_F\) is diagonal and
\[
        (A_FX_WA_F)_{vv}=|N_F(v)\cap W|.
\]
\end{lemma}

\begin{proof}
The \((u,v)\)-entry of \(A_FX_WA_F\) counts vertices \(w\in W\) such that
\(u\sim_F w\sim_F v\).  If \(u\ne v\), no leaf \(w\) can be adjacent to both \(u\) and
\(v\).  Hence all off-diagonal entries vanish.  For \(u=v\), the same count is
exactly the number of leaves in \(W\) adjacent to \(v\).
\end{proof}

The lemma is the precise point where forests differ from general graphs.  For an
arbitrary set \(W\), the matrix \(AX_WA\) records common \(W\)-neighbors and need not
be diagonal.

\subsection{Guarded pruning of a forest}
\label{subsec:guarded-pruning}

A lower rooted type is defined recursively during the pruning.  Initially there is
one active type.  When leaves of types \(\rho\) are peeled from a surviving vertex
of previous type \(\sigma\), the new type records \(\sigma\) together with the
multiplicities of the peeled child types.  Equivalently, it is the rooted
isomorphism type of the branch below that vertex, away from the eventual center,
as revealed by the current stage.

We keep the full vertex set throughout the pruning process.  At stage \(t\) we
maintain:
\begin{itemize}[leftmargin=2em]
    \item the adjacency matrix \(A_t\) and degree matrix \(D_t\) of the current active
    forest together with previously peeled isolated vertices;
    \item diagonal idempotents \(P_{t,\rho}\) selecting active vertices whose
    current lower rooted type is \(\rho\).
\end{itemize}
Let
\[
        R_t=\sum_\rho P_{t,\rho}
\]
be the active-vertex selector.  Initially
\[
        A_0=A(F),\qquad D_0=D(F),\qquad P_{0,\bullet}=I.
\]
All matrices constructed below lie in \(\Aalg(F)\).

For the selector steps used in the pruning, fix
\[
        \Sel_k^{(N)}(z)=\prod_{0\le r\le N,\ r\ne k}\frac{z-r}{k-r}.
\]
If a diagonal matrix \(Z\) has entries in \(\{0,\ldots,N\}\), then
\(\Sel_k^{(N)}(Z)\) is the diagonal idempotent of the \(k\)-level set of \(Z\).  We
suppress \(N\) when the range is clear.

The active degree-one selector is
\[
        L_t=R_t\Sel_1(D_t).
\]
By Lemma~\ref{lem:leaf-diagonal},
\[
        h_t=A_tL_tA_t
\]
is diagonal.  For an active leaf \(v\), the value \((h_t)_{vv}\) is \(1\) exactly when
its unique active neighbor is also an active leaf, i.e.\ when \(v\) lies in an
active \(K_2\) component.  Hence
\[
        C_t=L_t\Sel_1(h_t)
\]
selects the terminal active \(K_2\) components.  These vertices should not be
peeled; they are center-edge vertices.  The peelable type-\(\rho\) leaves are
selected by
\[
        \eta_{t,\rho}=P_{t,\rho}L_t(I-C_t),
        \qquad
        \eta_t=\sum_\rho \eta_{t,\rho}.
\]
If \(\eta_t=0\), the active forest is a disjoint union of \(K_1\) and \(K_2\)
components and the process stops.

For each lower type \(\rho\), set
\[
        g_{t,\rho}=A_t\eta_{t,\rho}A_t.
\]
Again by Lemma~\ref{lem:leaf-diagonal}, \(g_{t,\rho}\) is diagonal and
\[
        (g_{t,\rho})_{vv}=|N_{F_t}(v)\cap L_{t,\rho}^{\rm peel}|,
\]
where \(L_{t,\rho}^{\rm peel}\) is the set of peelable leaves of type \(\rho\).
Therefore \(\Sel_k(g_{t,\rho})\) selects the vertices adjacent to exactly \(k\)
currently peeled leaves of type \(\rho\).

The forest updates are
\[
        A_{t+1}=(I-\eta_t)A_t(I-\eta_t),
\]
\[
        D_{t+1}=D_t-\eta_t-\sum_\rho g_{t,\rho}.
\]
The first formula deletes all edges incident with the peeled leaves.  The second
subtracts one from each peeled leaf and subtracts, from every surviving neighbor,
the number of incident peeled leaves.

If a surviving vertex had old lower type \(\sigma\) and is adjacent to
\(k_\rho\) peeled leaves of type \(\rho\) for each \(\rho\), its new lower type is the
rooted type
\[
        \mu=(\sigma;(k_\rho)_\rho).
\]
The corresponding active selector is
\[
        P_{t+1,\mu}
        =
        P_{t,\sigma}(I-\eta_t)
        \prod_\rho \Sel_{k_\rho}(g_{t,\rho}),
\]
where only the finitely many types present at stage \(t\) occur in the product.
Thus every lower rooted type selector created by the pruning process is an
explicit diagonal idempotent in \(\Aalg(F)\).

\subsection{Center selectors}
\label{subsec:center-selectors}

Let \(s\) be the first stage with \(\eta_s=0\).  The active forest consists of
isolated active vertices and active \(K_2\) components.  Put
\[
        Z_s=R_s\Sel_0(D_s),
        \qquad
        L_s=R_s\Sel_1(D_s),
        \qquad
        h_s=A_sL_sA_s,
        \qquad
        C_s=L_s\Sel_1(h_s).
\]
Thus \(Z_s\) selects active center vertices of \(K_1\) components and \(C_s\) selects
active center-edge vertices of \(K_2\) components.

For each rooted type \(\rho\), the vector
\[
        z_\rho=Z_sP_{s,\rho}\one
\]
selects all single-center components whose center has lower rooted type \(\rho\).
These are center orbits of the forest.

For center edges, define the diagonal selector
\[
        E_{\rho\mid\sigma}
        =
        P_{s,\rho}C_s\Sel_1(A_sP_{s,\sigma}C_sA_s).
\]
It selects active \(K_2\) endpoints of type \(\rho\) whose center-edge neighbor has
type \(\sigma\).  If \(\rho\ne\sigma\), then
\[
        E_{\rho\mid\sigma}\one,
        \qquad
        E_{\sigma\mid\rho}\one
\]
are the two center orbits in all components with center types \(\{\rho,\sigma\}\).
If \(\rho=\sigma\), then
\[
        E_{\rho\mid\rho}\one
\]
selects both endpoints of each symmetric center edge, and these endpoints form a
single orbit.  All these seed orbit indicators belong to \(M_F\).

\subsection{Reverse propagation of orbit indicators}
\label{subsec:reverse-propagation}

Suppose that at stage \(t+1\) an active parent orbit \(O\) has already been
constructed, so that \(\one_O\in M_F\).  Let \(\eta_{t,\rho}\) be the selector of
peelable leaves of type \(\rho\) removed at stage \(t\).  Then
\[
        q=\eta_{t,\rho}A_t\one_O
\]
belongs to \(M_F\).  For a peeled leaf \(x\) of type \(\rho\), the coordinate
\((A_t\one_O)_x\) is \(1\) precisely when the unique active neighbor of \(x\) lies in
\(O\), and is \(0\) otherwise.  Hence
\[
        q=\one_{\{x\in L_{t,\rho}^{\rm peel}:p(x)\in O\}},
\]
where \(p(x)\) is the active parent of \(x\) at the time of peeling.

\subsection{Orbit classification in a forest}

Root every tree component at its center vertex, or at its center edge.  For a
non-center vertex \(v\), let \(p(v)\) be its parent and let \(\lambda(v)\) be the lower
rooted type of the component hanging below \(v\) away from the center.

\begin{lemma}[Rooted orbit classification]
\label{lem:rooted-orbit-classification}
Two non-center vertices \(v,w\) of a forest \(F\) are in the same automorphism orbit
if and only if
\[
        p(v)\text{ and }p(w)\text{ are in the same orbit}
        \quad\text{and}\quad
        \lambda(v)=\lambda(w).
\]
The center vertices are classified by the seed orbits described in
Subsection~\ref{subsec:center-selectors}.
\end{lemma}

\begin{proof}
Automorphisms preserve component centers, the parent relation, and rooted lower
types.  This gives the necessity.

Conversely, assume \(p(v)\) and \(p(w)\) are in the same orbit and
\(\lambda(v)=\lambda(w)\).  Choose an automorphism sending \(p(v)\) to \(p(w)\).
Inside the rooted branch below the parent, replace the image of the branch rooted
at \(v\) by a rooted isomorphism from the branch at \(v\) to the branch at \(w\).
Because equal lower rooted types occur as freely permutable rooted branches at a
fixed parent type, this local replacement extends to an automorphism of the whole
forest.  The center assertion follows from the usual classification of tree
centers: a component has either one center vertex or one center edge, and center
edges with equal endpoint rooted types have the endpoint-swap automorphism.
\end{proof}

\begin{theorem}[Forest orbit module theorem]
\label{thm:forest-module}
For every forest \(F\),
\[
        \Aalg(F)\one=U_F.
\]
In particular, for every tree \(T\),
\[
        \Aalg(T)\one=U_T.
\]
\end{theorem}

\begin{proof}
The inclusion \(\Aalg(F)\one\subseteq U_F\) follows from the observation in Section~\ref{sec:algebra} that automorphisms commute with \(A,D\) and fix \(\one\).
For the reverse inclusion, the guarded pruning process constructs all lower
rooted type selectors inside \(\Aalg(F)\).  At the terminal stage,
Subsection~\ref{subsec:center-selectors} gives all center orbit indicators as vectors in
\(M_F\).  Applying the reverse propagation step of
Subsection~\ref{subsec:reverse-propagation} from the terminal stage back to stage \(0\)
constructs every non-center orbit indicator.  By
Lemma~\ref{lem:rooted-orbit-classification}, the constructed sets are exactly the
automorphism orbits of the forest.  Thus \(U_F\subseteq M_F\).
\end{proof}


\section{Tree quotient full-matrix rigidity}
\label{sec:tree-quotient-full-matrix}

The orbit-module theorem says that the vector \(\one\) is cyclic for the action of
\(\langle A(T),D(T)\rangle\) on the orbit module.  For connected trees a stronger
statement holds: the entire quotient representation is the full matrix algebra.


Let \(T\) be a connected tree and let \(O_1,\ldots,O_m\) be the orbits of
\(\Aut(T)\).  Put
\[
        S=(\one_{O_1},\ldots,\one_{O_m}).
\]
The orbit partition is equitable, and the cells are degree-homogeneous, so there
are unique matrices \(B_T\) and \(\Delta_T\) with
\[
        A(T)S=SB_T,\qquad D(T)S=S\Delta_T,
\]
where \(\Delta_T\) is diagonal.  Define the quotient algebra
\[
        Q_T:=\langle I_m,B_T,\Delta_T\rangle\subseteq M_m(\R).
\]

\begin{theorem}[Tree quotient full-matrix theorem]
\label{thm:tree-quotient-full-matrix}
For every finite connected tree \(T\),
\[
        Q_T=M_m(\R).
\]
Equivalently,
\[
        \langle A(T),D(T)\rangle|_{U_T}=\End(U_T).
\]
\end{theorem}

We give a constructive proof.  The proof produces every coordinate idempotent on
the orbit quotient and then every matrix unit.  Thus no appeal to the complex
Burnside theorem, or to real Wedderburn classification, is needed.


\begin{lemma}[Rooted quotient geometry]
\label{lem:rooted-quotient-geometry}
Root a connected tree at its center vertex or at its center edge.  The support of
its orbit quotient is a tree rooted at a vertex or at an edge, except that a
bicentral edge whose two endpoint rooted branches are isomorphic is represented
by a loop at the single root orbit.  After deleting this possible loop, every
non-root quotient vertex has a unique parent.  Moreover, if quotient vertices
$i,i'$ are children of the same parent quotient vertex and have the same lower
rooted type, then $i=i'$.
\end{lemma}

\begin{proof}
The parent relation and lower rooted type are invariant under automorphisms.
Lemma~\ref{lem:rooted-orbit-classification} gives, on non-center orbits,
\[
        i=i'
        \Longleftrightarrow
        p(i)=p(i')\quad\text{and}\quad \lambda(i)=\lambda(i') .
\]
Thus every non-root quotient vertex has a unique parent, and there is at most one
child of a fixed lower type below a fixed parent.  At the center there is either
one root orbit, two root orbits joined by the center edge, or one root orbit with
a loop representing a symmetric center edge.  This proves the asserted quotient
geometry.
\end{proof}

\begin{lemma}[Quotient idempotent extraction]
\label{lem:quotient-idempotents}
For a connected tree \(T\), every coordinate idempotent \(E_i\in M_m(\R)\) of the
orbit quotient belongs to \(Q_T=\langle I_m,B_T,\Delta_T\rangle\).
\end{lemma}

\begin{proof}
Apply the guarded pruning formulas of Section~\ref{sec:trees} to the quotient.
At every stage
\[
        B_{t+1}=(I-\eta_t)B_t(I-\eta_t),\qquad
        \Delta_{t+1}=\Delta_t-\eta_t-\sum_\rho B_t\eta_{t,\rho}B_t,
\]
so all matrices constructed remain in \(Q_T\).  If \(X\) selects current quotient
leaves of lower type \(\rho\), Lemma~\ref{lem:rooted-quotient-geometry} gives
\[
        B_tXB_t\in Q_T\cap \operatorname{Diag}_m(\R),
\]
and its diagonal entries are integers in \(\{0,\ldots,|V(T)|\}\).  Hence the
idempotents \(\Sel_k(B_tXB_t)\) used in the pruning belong to \(Q_T\).  Thus the
same formulas as in Section~\ref{sec:trees} produce the diagonal idempotents of
the terminal center orbit or orbits.

Now propagate away from the center.  If \(E_p\in Q_T\) is the idempotent of a
known parent quotient vertex and \(X\) selects children of a fixed lower type,
then
\[
        Z=XB_tE_pB_tX\in Q_T\cap \operatorname{Diag}_m(\R).
\]
By Lemma~\ref{lem:rooted-quotient-geometry}, the support of \(Z\) is either empty
or the single desired child quotient vertex.  Therefore
\[
        \sum_{k=1}^{|V(T)|}\Sel_k(Z)
\]
is that child's coordinate idempotent.  Induction over the rooted quotient tree
gives all \(E_i\).
\end{proof}

\begin{proof}[Proof of Theorem~\ref{thm:tree-quotient-full-matrix}]
By Lemma~\ref{lem:quotient-idempotents}, each coordinate idempotent \(E_i\) belongs to
\(Q_T\).  If quotient vertices \(i,j\) are adjacent, then
\[
        E_iB_TE_j=(B_T)_{ij}E_{ij},
\]
with \((B_T)_{ij}>0\), so the directed matrix unit \(E_{ij}\) lies in \(Q_T\).
The support quotient is connected after ignoring the possible root loop.  For any
ordered pair \((i,j)\), choose a quotient path
\[
        i=i_0,i_1,\ldots,i_r=j.
\]
Then the product
\[
        E_{i_0i_1}E_{i_1i_2}\cdots E_{i_{r-1}i_r}=E_{ij}
\]
belongs to \(Q_T\).  Thus every matrix unit belongs to \(Q_T\), and therefore
\(Q_T=M_m(\R)\).
\end{proof}

In particular, for a connected tree \(T\), the simultaneous commutant of the
quotient operators is scalar:
\[
        \{X\in M_m(\R):XB_T=B_TX,\ X\Delta_T=\Delta_TX\}=\R I_m.
\]
This conclusion is tree-specific.  For general graphs the algebra induced by
\(A\) and \(D\) on the orbit quotient can be much smaller than a full matrix
algebra; the full-matrix conclusion for trees is a consequence of the rooted
leaf-pruning structure of tree orbit quotients.

\section{Tree reconstruction from principal moments}
\label{sec:tree-reconstruction}

\subsection{Orbit quotients and rank-one moments}

Let \(O_1,\ldots,O_m\) be the vertex orbits of a graph \(G\), and set
\(u_i=\one_{O_i}\).  The orbit partition is equitable.  Define
\[
        s_i=u_i^Tu_i=|O_i|,
        \qquad
        q_{ij}=u_i^TAu_j,
        \qquad
        b_{ij}=\frac{q_{ij}}{s_i}.
\]
Then \(b_{ij}\) is the number of neighbors in \(O_j\) of any fixed vertex in \(O_i\).
We call the data
\[
        (s_i,b_{ij})_{1\le i,j\le m}
\]
the weighted orbit quotient.  Diagonal entries are allowed; for a tree they occur
only when the two endpoints of the center edge lie in the same orbit.

If \(u_i=p_i(A,D)\one\), then these quotient data are rank-one moments:
\[
        s_i=\one^Tp_i(A,D)^Tp_i(A,D)\one,
\]
\[
        q_{ij}=\one^Tp_i(A,D)^TAp_j(A,D)\one.
\]
Thus the principal block together with the distinguished adjacency operator
recovers the weighted orbit quotient once orbit selectors are known.


\subsection{Orbit-quotient rigidity of trees}

\begin{theorem}[Orbit-quotient rigidity]
\label{thm:tree-quotient-rigidity}
A tree is determined up to isomorphism by its weighted orbit quotient.
\end{theorem}

\begin{proof}
Let \(Q=(s_i,b_{ij})\) be the weighted orbit quotient of a tree.  For an active
set \(S\subseteq\{1,\ldots,m\}\), put
\[
        d_i(S)=\sum_{j\in S} b_{ij}.
\]
Starting with \(S_0=\{1,\ldots,m\}\), delete all quotient vertices with
\(d_i(S_t)=1\), except when the remaining active quotient is one of the two
center-edge cases
\[
        S=\{c\},\quad s_c=2,\quad b_{cc}=1,
        \qquad\text{or}\qquad
        S=\{c,c'\},\quad b_{cc'}=b_{c'c}=1.
\]
This is exactly the image of simultaneous leaf stripping on the original tree,
because an active vertex in orbit \(O_i\) has active degree \(d_i(S_t)\).  Hence
the process determines the center type.  If \(i\) is deleted at stage \(t\), its
unique active neighbor is the unique \(p(i)\in S_t\) with \(b_{i,p(i)}=1\); this
is the parent orbit of \(i\).

Define rooted types from the leaves inward.  If \(i\) is a deleted leaf orbit,
set \(\tau_i=\bullet\).  After all children of \(i\) have been assigned types,
set
\[
        \tau_i=\bullet\Bigl(\{\tau_j^{\,b_{ij}}:p(j)=i\}\Bigr),
\]
where the notation records the multiset of rooted child types attached to one
vertex of orbit \(O_i\).  This recursion is forced by \(Q\).  At the terminal
core, either a single center type \(\tau_c\) remains, or two rooted types are
glued across a center edge, or two copies of the same rooted type are glued
across a symmetric center edge represented by a loop \(b_{cc}=1\).  Thus \(Q\)
constructs a unique tree up to isomorphism.
\end{proof}

\begin{lemma}[Cyclic realization of the principal moment functional]
\label{lem:gns-principal-realization}
The functional
\[
        \mu_G(w)=\one^T w(A_G,D_G)\one,
        \qquad w\in\mathbb R\langle x,y\rangle,
\]
determines the pointed finite-dimensional cyclic representation
\[
        (M_G,A_G|_{M_G},D_G|_{M_G},\one,\langle\cdot,\cdot\rangle).
\]
More precisely, it determines it up to an isometry carrying $\one$ to $\one$ and
intertwining the two represented generators.
\end{lemma}

\begin{proof}
Form the pre-Hilbert space on $\mathbb R\langle x,y\rangle$ with bilinear
form
\[
        \langle p,q\rangle_\mu=\mu_G(p^{\rm rev}q).
\]
Quotient by the null space $\{p:\mu_G(p^{\rm rev}p)=0\}$ and complete; the result
is already finite-dimensional because it is represented by the vectors
$p(A_G,D_G)\one$ in $M_G$.  Left multiplication by $x$ and $y$ gives the induced
operators corresponding to $A_G$ and $D_G$, and the class of the empty word is
the cyclic vector.  This is the usual finite GNS construction for the real
symmetric moment functional.
\end{proof}

\begin{lemma}[Intrinsic recovery of the weighted quotient]
\label{lem:intrinsic-weighted-quotient}
For a tree $T$, the pointed cyclic representation
\[
        (M_T,A_T|_{M_T},D_T|_{M_T},\one,\langle\cdot,\cdot\rangle)
\]
determines the weighted orbit quotient of $T$, up to simultaneous relabeling of
its quotient vertices.
\end{lemma}

\begin{proof}
By Theorem~\ref{thm:forest-module}, $M_T=U_T$.  Let
$u_i=\one_{O_i}$ be the orbit indicators.  In the basis $(u_i)$, the represented
operators $A_T|_{U_T}$ and $D_T|_{U_T}$ have matrices $B_T$ and $\Delta_T$,
respectively.  The idempotent-extraction procedure of
Lemma~\ref{lem:quotient-idempotents} is expressed entirely by algebraic
operations in the represented pair, together with the fixed level-selector
polynomials $\Sel_k^{(N)}$.  Therefore the same procedure, transported through
any isometric intertwiner of pointed cyclic representations, recovers the
one-dimensional coordinate projections $P_i$ onto the lines $\R u_i$, up to
permutation of the indices.

Since the cyclic vector is $\one=\sum_i u_i$, we have $u_i=P_i\one$.  Hence the
representation determines
\[
        s_i=\|P_i\one\|^2,
        \qquad
        q_{ij}=\langle P_i\one,\, A_T P_j\one\rangle .
\]
The weighted quotient entries are $b_{ij}=q_{ij}/s_i$.  Thus the weighted orbit
quotient is determined up to relabeling.
\end{proof}

\begin{theorem}[Principal AD moments determine trees]
\label{thm:principal-moments-determine-trees}
Let $T$ and $T'$ be finite trees.  Suppose that
\[
        \mu_T(w)=\mu_{T'}(w)
        \qquad\text{for every }w\in\mathbb R\langle x,y\rangle,
\]
where $\mu_G(w)=\one^T w(A_G,D_G)\one$.  Then $T\cong T'$.
\end{theorem}

\begin{proof}
By Lemma~\ref{lem:gns-principal-realization}, equality of the moment functionals gives
an isometric isomorphism of the pointed cyclic representations generated by
$\one$.  Lemma~\ref{lem:intrinsic-weighted-quotient} then shows that the two trees have
the same weighted orbit quotient, up to relabeling.  By
Theorem~\ref{thm:tree-quotient-rigidity}, this quotient determines the tree.
\end{proof}

\begin{remark}
The abstract algebra \(\End(U_T)\) alone does not determine \(T\); it remembers only
\(\dim U_T\).  The distinguished \(A\)-moments are essential.
\end{remark}



It is useful to compare the theorem with a familiar one-parameter refinement.
Wang--Li--Lu--Xu introduced the generalized characteristic polynomial
\cite{WangLiLuXu2011}
\[
        \phi_G(\lambda,t)=\det(\lambda I-(A_G-tD_G)).
\]
This pencil is a natural slice of the adjacency-degree calculus.  McKay's
coalescence construction already gives families of trees with the same
$A-tD$ spectrum for every value of $t$; related limitations of immanantal
polynomials for trees were proved by Botti and Merris~\cite{BottiMerris1993}.
For a small concrete witness, let $T_1,T_2$ be the trees on
$\{0,1,\ldots,10\}$ with edge sets
\[
\begin{aligned}
E(T_1)=\{&01,02,05,06,17,23,24,78,7\,10,89\},\\
E(T_2)=\{&01,02,15,1\,10,23,24,56,58,59,67\}.
\end{aligned}
\]
They have the same degree sequence but different branch structure at the unique
vertex of degree four.  A direct matching expansion gives
\[
        \phi_{T_1}(\lambda,t)=\phi_{T_2}(\lambda,t),
\]
and the common polynomial is recorded in Appendix~\ref{app:wang-pencil}.  The
principal moment theorem therefore uses more than this one-parameter pencil.


\section{Finite certificates and algorithmic extraction}\label{sec:algorithmic-extraction}
\subsection{Finite graph-dependent witnesses}

\begin{proposition}[Finite graph-dependent witnesses]
\label{prop:finite-witnesses}
Let $T$ be a tree and let $m$ be the number of automorphism orbits of $T$.  There
exist words $w_1,\ldots,w_m\in\R\langle x,y\rangle$ such that the finite matrices
\[
        \bigl(\mu_T(w_i^{\rm rev}w_j)\bigr),
        \qquad
        \bigl(\mu_T(w_i^{\rm rev}xw_j)\bigr),
        \qquad
        \bigl(\mu_T(w_i^{\rm rev}yw_j)\bigr)
\]
determine the weighted orbit quotient of $T$, and hence determine $T$.
The words may be chosen by the greedy AD-Krylov basis construction and are
therefore graph-dependent.
\end{proposition}

\begin{proof}
Choose words $w_i$ so that $w_i(A,D)\one$ is a basis of $M_T=U_T$.  The first
matrix is the Gram matrix of this basis.  The second and third give the bilinear
forms of $A$ and $D$ on the same basis; multiplying by the inverse Gram matrix
gives the matrices of the quotient operators.  The reconstruction argument of Section~\ref{sec:tree-reconstruction} then recovers the orbit
idempotents and weighted orbit quotient.  A deterministic shortlex greedy
procedure gives such a basis using at most $m$ accepted words.
\end{proof}

\begin{remark}
The proposition is intentionally graph-dependent.  It does not assert the
existence of a small universal word family, depending only on $n$, that works for
all $n$-vertex trees.  Such a finite universal McKay-word collapse is a natural
separate problem.
\end{remark}


\subsection{Algorithmic extraction for trees}
\label{subsec:algorithmic-tree-extraction}

The constructive proof gives an algorithmic corollary.  This is not meant as a
new fastest tree-isomorphism algorithm: classical combinatorial tree isomorphism
is already linear-time.  The point is that the principal AD block supplies a
linear-algebraic certificate for the same canonical quotient data.

\begin{proposition}[Linear-time tree canonization and AD certificates]
\label{prop:linear-time-tree-canonization}
There is a linear-time algorithm which, given a tree $T$, computes its
automorphism-orbit partition, its weighted orbit quotient
\[
        (s_i,b_{ij})_{1\le i,j\le m},
\]
and a canonical tree code.  Moreover, in an algebraic straight-line model that
allows the level-selector polynomials $\Sel_k^{(N)}$ as gates, the same run
records graph-dependent AD expressions whose evaluations are the orbit
indicators of $T$.
\end{proposition}

\begin{proof}
Run the guarded pruning process of Section~\ref{sec:trees}.  With bucketed rooted-type
identifiers, each edge and vertex is processed a bounded number of times, so the
usual AHU/McKay leaf-stripping implementation is linear.  The construction in
Subsections~\ref{subsec:guarded-pruning}, \ref{subsec:center-selectors}, and~\ref{subsec:reverse-propagation}
records the same steps as AD expressions: leaf selectors, parent-count
diagonals, center selectors, and reverse propagation all lie in the algebra
generated by $A$ and $D$.  Once the orbit indicators are obtained, a single scan
of the edges gives the weighted quotient entries $b_{ij}$ and the cell sizes
$s_i$.  Theorem~\ref{thm:tree-quotient-rigidity} then gives a canonical code by pruning
the weighted orbit quotient from the leaves toward the center.
\end{proof}

\begin{remark}
The straight-line expressions in Proposition~\ref{prop:linear-time-tree-canonization} are
allowed to depend on the input tree.  If all selector polynomials are expanded as
ordinary noncommutative expressions, no linear bound on expression length is
claimed here.  The problem of finding a small universal family of AD words or
decorated-caterpillar queries, depending only on $n$, is a separate finite-basis
problem.
\end{remark}



\section{Principal AD-rigidity and color refinement}
\label{sec:caterpillar-rigidity}

The tree theorem suggests a graph class.  We call two graphs $G,H$ \emph{principal
equivalent}, written $G\equiv_{\rm prin}H$, if
\[
        \one^Tw(A_G,D_G)\one=\one^Tw(A_H,D_H)\one
        \qquad\text{for every }w\in\R\langle x,y\rangle .
\]
A graph is \emph{principal AD-rigid} if its principal equivalence class contains
no non-isomorphic graph.  This is a structural notion: it is the unique-realization
class for the principal cyclic noncommutative AD moment problem.

The module $M_G$ also induces a vertex partition, the \emph{principal profile
partition}: vertices $u,v$ have the same profile if $f(u)=f(v)$ for every
$f\in M_G$.  These profile classes are generally coarser than the stable
color-refinement cells and need not be equitable.  Their role below is to locate
switching directions invisible to the principal moments.

For nonnegative integers $a_0,\ldots,a_\ell$, let $C(a_0,\ldots,a_\ell)$ be the
caterpillar obtained from the spine $0-1-\cdots-\ell$ by attaching $a_i$ pendant
leaves to the $i$-th spine vertex.

\begin{proposition}[Decorated-caterpillar translation]
\label{prop:caterpillar-translation}
For every graph $G$,
\[
        \homG(C(a_0,\ldots,a_\ell),G)
        =
        \one^TD^{a_0}AD^{a_1}A\cdots AD^{a_\ell}\one .
\]
Thus monomial principal AD moments are precisely homomorphism counts from
degree-decorated caterpillars, and polynomial moments are their linear
combinations.
\end{proposition}

\begin{proof}
Expanding the matrix coefficient gives
\[
        \sum_{v_0\sim v_1\sim\cdots\sim v_\ell}
        \prod_{i=0}^{\ell}\deg(v_i)^{a_i}.
\]
For a fixed image of the spine, the $a_i$ leaves attached at the $i$-th spine
vertex may be mapped independently to any neighbor of $v_i$, giving
$\deg(v_i)^{a_i}$ choices.  Summing over all spine walks proves the identity.
\end{proof}

This identifies the pattern class behind principal AD-rigidity.  Color
refinement, by the theorem of Dell--Grohe--Rattan, is equivalence under all tree
homomorphism counts~\cite{DellGroheRattan2018}.  Principal AD moments see only the
one-spine, degree-decorated caterpillar subfamily.  Therefore color-refinement
equivalence implies principal equivalence.  In particular, every principal
AD-rigid graph is amenable to color refinement in the sense of
Arvind--Kobler--Rattan--Verbitsky~\cite{ArvindKoblerRattanVerbitsky2017}.

Combining this with the compact-graph hierarchy of
\citet{ArvindKoblerRattanVerbitsky2017}, we obtain the placement
\[
        \{\text{principal AD-rigid graphs}\}
        \subseteq \{\text{amenable graphs}\}
        \subseteq \{\text{compact graphs}\}
        \subseteq \{\text{refinable graphs}\}.
\]
Only the first inclusion is specific to the principal moment functional.  The
others belong to the existing color-refinement and compactness theory.  None of
the converses is asserted here.  Compactness is already too broad for the module
question, as noted in the introduction; the 10-vertex switch below is stronger,
being amenable, hence compact and refinable, but not principal AD-rigid.

The distinction is structural.  Rooted tree messages are generated by
\[
        f\mapsto Af,
        \qquad
        (f,g)\mapsto f\circ g,
\]
where $\circ$ denotes coordinatewise product; the product is the algebraic form
of branching.  Principal AD messages are generated only by
\[
        f\mapsto Af,
        \qquad
        f\mapsto Df=(A\one)\circ f.
\]
Thus principal AD-rigidity is rigidity under a one-spine partial tree-Hankel
system, whereas amenability is rigidity under the full rooted-tree branching
system.

The first obstruction to principal AD-rigidity is an integral switching invisible
to the principal module.

\begin{definition}[Principal-invisible integral switching]
Let $G$ have adjacency matrix $A$ and principal module $M_G$.  A
\emph{principal-invisible integral switching} is a nonzero symmetric integral
zero-diagonal matrix $E$ such that
\[
        EM_G=0
\]
and $A+E$ is again the adjacency matrix of a simple graph.
\end{definition}

\begin{lemma}[Invisible switching preserves principal moments]
\label{lem:invisible-switching}
Let $H$ be the graph with adjacency matrix $A_H=A_G+E$, where $E$ is a
principal-invisible integral switching of $G$.  Then $G\equiv_{\rm prin}H$.
\end{lemma}

\begin{proof}
Since $\one\in M_G$, the condition $EM_G=0$ implies $E\one=0$, so $D_H=D_G=:D$.
On $M_G$ the operators $A_H=A_G+E$ and $A_G$ agree.  By induction on word length,
$w(A_H,D)\one=w(A_G,D)\one$ for every word $w$, and hence all principal moments
are equal.
\end{proof}

In particular, if $a,a'$ have the same principal profile and $b,b'$ have the same
principal profile, then a legal switch
\[
        ab,\\ a'b' \longleftrightarrow ab',\\ a'b
\]
is invisible to principal moments.  Indeed the corresponding perturbation is
\[
        (e_a-e_{a'})(e_{b'}-e_b)^T+(e_{b'}-e_b)(e_a-e_{a'})^T.
\]



The graph $G_{10}$ of Example~\ref{ex:n10-counterexample} is color-refinement-discrete and hence amenable, but
$G_{10}\equiv_{\rm prin}H_{10}$ and $G_{10}\not\cong H_{10}$.  Appendix~\ref{app:n10-certificate}
gives a compact certificate independent of the enumeration code.  With the
displayed labeling, the graph6 strings are
\[
        G_{10}=\texttt{I?hcivILW},
        \qquad
        H_{10}=\texttt{I?YcivILW}.
\]
The complete enumeration records an isomorphic representative of the first graph
as \texttt{I?qbDrs\{\_}.
Thus principal AD-rigid graphs form a proper subclass of the amenable graphs.



The example is representative of the first failures.  Complete enumeration up to
nine vertices gives a stronger small-order coincidence: the principal equivalence
partition and the color-refinement equivalence partition are identical.  For
$n=8$ there are $12{,}346$ graphs, $12{,}095$ equivalence classes, and $11{,}892$
singletons; for $n=9$ there are $274{,}668$ graphs, $271{,}941$ equivalence
classes, and $269{,}750$ singletons.  The verification compares exact canonical
color-refinement certificates with complete principal Krylov signatures; hash
values are used only as compact labels.  At $n=10$ the first failures appear.  A
full targeted enumeration of the $12{,}005{,}168$ unlabeled graphs on ten vertices
found 107 nontrivial integral principal-fiber witnesses and 101 direct graph
certificates; all 101 were independently rechecked to have equal principal
signatures and different color-refinement signatures.

All finite enumerations reported here were carried out with exact integer or
rational arithmetic.  For each graph, the principal cyclic representation was
closed under left multiplication by $A$ and $D$ using fraction-free rank tests,
and the resulting finite Gram and represented-operator data were canonically
reduced before comparison; hash values were used only as labels for these exact
certificates.  The explicit ten-vertex witness in Example~\ref{ex:n10-counterexample}
is certified directly in Appendix~\ref{app:n10-certificate}.

The structure of these certificates is highly regular.  All 101 direct source
graphs are color-refinement discrete but have nontrivial principal profile
classes.  The hidden perturbations remove and add the same number of undirected
edges; the counts are
\[
        (2,2):86,\qquad (4,4):12,
        \qquad (6,6):2,
        \qquad (8,8):1.
\]
For the 86 single $2$-switches, the endpoint degree-pair multiset is preserved;
the most common patterns are two $(3,4)$-edges, two $(4,5)$-edges, and two
$(3,5)$-edges.  In all 101 direct certificates, the multiset of endpoint
principal-profile pairs is preserved.  The principal profile size patterns of
the source graphs are
\[
\begin{array}{c|c}
\text{profile sizes} & \text{number of certificates}\\ \hline
(1,1,2,2,2,2) & 84\\
(1,1,2,2,4) & 7\\
(1,1,1,1,2,2,2) & 6\\
(1,1,1,1,1,1,2,2) & 4.
\end{array}
\]
Thus the observed gap from amenability arises from graph-realizable integral
switchings in directions left unresolved by the principal cyclic module.

The resulting classification problem is to characterize principal AD-rigid graphs
intrinsically.  A particularly concrete form is: determine whether every amenable
graph that fails principal AD-rigidity admits, possibly after passing to the
stable quotient, a nontrivial graph-realizable switching invisible to the
principal cyclic module.  A second problem, already suggested by the constructive
tree proof, is to find small universal finite families of decorated-caterpillar
moments which determine all $n$-vertex forests.



\appendix

\section{Finite certificates for the ten-vertex switch}
\label{app:n10-certificate}

This appendix records a compact, checkable certificate for Example~\ref{ex:n10-counterexample}.  The certificate is independent of the enumeration code: one checks the displayed closure identities for $M_{G_{10}}$, the repeated principal profiles used by the switch, the color-refinement refinement sequence, and the final triangle-count separation.
Let $A$ be the adjacency matrix of $G_{10}$ and set
\[
        K=\bigl[\one,\ A\one,\ A^2\one,\ DA\one,\ A^3\one\bigr].
\]
Explicitly, with rows ordered by vertices $0,
1,
\ldots,
9$,
\[
K=
\begin{pmatrix}
1&3&14&9&56\\
1&3&14&9&56\\
1&3&13&9&52\\
1&3&15&9&60\\
1&4&16&16&67\\
1&4&16&16&67\\
1&5&20&25&85\\
1&5&20&25&85\\
1&5&20&25&84\\
1&5&20&25&84
\end{pmatrix}.
\]
It has rank $5$.  Moreover the following identities hold:
\[
        AK=KC_A,
        \qquad
        DK=KC_D,
\]
where
\[
C_A=
\begin{pmatrix}
0&0&0&-36&17\\
1&0&0&-2&-14\\
0&1&0&9&2\\
0&0&0&-2&1\\
0&0&1&0&4
\end{pmatrix},
\]
\[
C_D=
\begin{pmatrix}
0&0&0&60&45\\
1&0&-12&-47&-72\\
0&0&3&0&-8\\
0&1&4&12&19\\
0&0&0&0&5
\end{pmatrix}.
\]
Thus the five displayed columns span the full principal module $M_{G_{10}}$.
Since rows $0$ and $1$ of $K$ are equal, and rows $4$ and $5$ are equal, the
switch $04,15\leftrightarrow 05,14$ has the principal-profile form discussed before Example~\ref{ex:n10-counterexample}; equivalently, its perturbation matrix kills $M_{G_{10}}$.

For color refinement, the successive cell-size patterns for $G_{10}$ are
\[
(4,4,2),\quad (4,2,2,1,1),\quad (2,2,2,1,1,1,1),
\]
\[
(2,1,1,1,1,1,1,1,1),\quad (1,1,1,1,1,1,1,1,1,1).
\]
Thus the stable partition is discrete.  Finally,
\[
        \operatorname{tr}(A_{G_{10}}^3)=30,
        \qquad
        \operatorname{tr}(A_{H_{10}}^3)=42,
\]
so the graphs have different numbers of triangles and are not isomorphic.

\section{The common Wang pencil polynomial}
\label{app:wang-pencil}

For the two non-isomorphic trees displayed in Section~\ref{sec:tree-reconstruction}, the common
polynomial is
\[
\begin{aligned}
(\lambda+t)^2\bigl(&
\lambda^9+18\lambda^8t+139\lambda^7t^2-10\lambda^7
+604\lambda^6t^3-132\lambda^6t\\
&+1627\lambda^5t^4-717\lambda^5t^2+31\lambda^5
+2818\lambda^4t^5-2076\lambda^4t^3+272\lambda^4t\\
&+3141\lambda^3t^6-3462\lambda^3t^4+910\lambda^3t^2-33\lambda^3\\
&+2176\lambda^2t^7-3330\lambda^2t^5+1452\lambda^2t^3-160\lambda^2t\\
&+852\lambda t^8-1715\lambda t^6+1107\lambda t^4-243\lambda t^2+10\lambda\\
&+144t^9-366t^7+324t^5-116t^3+14t
\bigr).
\end{aligned}
\]



\section*{Acknowledgments} This work is supported by Guangdong Province (No. 2023QN10X215), 2023 Shenzhen National Science Foundation (No. 20231128220938001), Shenzhen Science and Technology Program (No. JCYJ20241202130548062), the Natural Science Foundation of Shenzhen (No. JCYJ20230807142703006), and the Key Research Platforms and Projects of the Guangdong Provincial Department of Education (No.2023ZDZX1034). 






\section*{Declaration of generative AI and AI-assisted technologies in the manuscript preparation process}

During the preparation of this work, the authors used ChatGPT and Codex, for language editing, organization of exposition, proofreading, consistency checks, assistance with
bibliographic searches, verification of selected computations. The mathematical ideas, definitions, main framework, and proofs were developed by the first author. The authors reviewed, revised, and verified the AI-assisted output as needed and take full responsibility for the content of the article.





\bibliographystyle{cas-model2-names}
\bibliography{cas-refs-revised}


\end{document}
